\subsection{可解群的情形}

设 \(\mathbf{E}\) 是一个实向量空间，\(\mathbf{K}\) 是 \(\mathbf{E}\) 的一个凸子集。一个映射 \(u: \mathbf{K} \to \mathbf{K}\)，若对 \(x, y\)（属于 \(\mathbf{K}\)）以及对每个实数 \(t\)（属于 \([0, 1]\)），我们有

\[
u (t x + (1 - t) y) = t u (x) + (1 - t) u (y)\tag{1}
\]

就称为一个仿射变换。由关系式 (1) 我们推出

\[
u \left(\sum_ {i \in \mathrm{I}} t _ {i} x _ {i}\right) = \sum_ {i \in \mathrm{I}} t _ {i} u \left(x _ {i}\right)\tag{2}
\]

对每个有限集 I、K 中的点 \(x_{i}\) 以及正实数 \(t_i\)（满足 \(\sum_{i\in I}t_i = 1\)）成立。

设 u 与 v 是 K 上的两个仿射变换，则映射 \(u \circ v\) 是 K 上的一个仿射变换。若 \(v: E \to E\) 是一个满足 \(v(\mathbf{K}) \subset \mathbf{K}\) 的线性映射，则在 K 上与 v 相同的映射 \(u: K \to K\) 是一个仿射变换。

定理 1 (Markoff-Kakutani). --- 设 E 是域 R 上的一个 Hausdorff 局部凸向量空间，K 是 E 的一个非空紧凸子集。设 \(\Gamma\) 是 K 上的一组两两可交换的连续仿射变换。那么存在 K 中的一点 a，使得 \(u(a) = a\) 对所有 \(u \in \Gamma\) 成立。

对每个 \(u \in \Gamma\)，设 \(\mathbf{K}_u\) 是所有 \(x \in \mathbf{K}\) 中满足 \(u(x) = x\) 者所成的集。我们将证明 \(\mathbf{K}_u\) 非空。设 \(x\) 是 \(\mathbf{K}\) 的一个点；对每个整数 \(n \geqslant 1\)，设 \(x_n\) 是元素 \(\frac{1}{n} \sum_{i=0}^{n-1} u^i(x)\)（\(\mathbf{E}\) 中的元素）。由于 \(\mathbf{K}\) 凸且在 \(u\) 下稳定，诸点 \(x_n\) 属于 \(\mathbf{K}\)，且由于 \(\mathbf{K}\) 紧，存在一个极限点 \(a\)（序列 \((x_n)_{n \geqslant 1}\) 的极限点）。映射 \(y \mapsto u(y) - y\)（从 \(\mathbf{K}\) 到 \(\mathbf{E}\) 中的映射）是连续的，因而 \(u(a) - a\) 是序列 \((u(x_n) - x_n)_{n \geqslant 1}\) 的一个极限点。但我们有 \(u(x_n) - x_n = \frac{1}{n} (u^n(x) - x)\)。

由于 K 紧，从而有界 (III, 第 3 页, 命题 2)，序列 \((u^n(x) - x)_{n\geqslant 1}\) 有界；于是序列 \(\left(\frac{1}{n} (u^n (x) - x)\right)_{n\geqslant 1}\) 趋于 0 (III, 第 4 页, 命题 3)，且由于 E 是 Hausdorff 的，我们有 \(u(a) - a = 0\)。因此 \(a\in \mathbf{K}_u\)。

诸集 \(\mathbf{K}_u\) 中的每一个都是紧空间 \(\mathbf{K}\) 的闭凸子集，我们将证明交 \(\bigcap_{u\in \Gamma}\mathbf{K}_u\) 非空。为此只需

证明：对 \(n \geqslant 1\) 以及 \(u_{1}, \ldots, u_{n}\)（属于 \(\Gamma\)），集 \(K_{u_{1}} \cap \ldots \cap K_{u_{n}}\) 非空。n = 1 的情形已经考虑过，我们对 n 用归纳法论证。设 \(n \geqslant 2\)，并令 \(L = K_{u_{1}} \cap \ldots \cap K_{u_{n-1}}\)。由归纳假设，L 是 E 的一个非空紧凸子集。由于 \(u_{n}\) 与 \(u_{1}, \ldots, u_{n-1}\) 交换，我们有 \(u_{n}(L) \subset L\)。把证明的第一部分应用于 \(u_{n}\) 在 L 上诱导的仿射变换，我们断定存在 L 中的一点 a 使得 \(u_{n}(a) = a\)；于是 a 属于 \(K_{u_{1}} \cap \ldots \cap K_{u_{n}}\)，从而该集非空。

推论. --- 设 G 是 K 上连续仿射变换的一个可解群。那么存在 K 中的一点，它在 G 下不变。

根据可解群的定义 (A, I, § 6, No.~4)，存在一个有限递降序列 \((\mathbf{G}_i)_{0\leqslant i\leqslant n}\)（由 \(\mathbf{G}\) 的互不相同的子群组成），使得 \(\mathbf{G}_0 = \mathbf{G}\)，\(\mathbf{G}_n = \{e\}\)，且群 \(\mathbf{G}_{i - 1} / \mathbf{G}_i\) 对 \(1\leqslant i\leqslant n\) 是交换的。设 \(\mathbf{K}_i\) 表示 \(\mathbf{G}_i\) 在 \(\mathbf{K}\) 中的不动点集。那么 \(\mathbf{K}_n = \mathbf{K}\)。此外，对 \(1\leqslant i\leqslant n\)，\(\mathbf{G}_i\) 的每个元素在 \(\mathbf{K}_i\) 上诱导恒等变换；这样，我们就得到交换群 \(\mathbf{G}_{i - 1} / \mathbf{G}_i\) 在 \(\mathbf{K}\) 上的一个作用（若 \(\mathbf{K}_i\) 非空）；由定理 1 可知，不动点集 \(\mathbf{K}_{i - 1}\)（\(\mathbf{G}_{i - 1} / \mathbf{G}_i\) 在 \(\mathbf{K}_i\) 中的不动点集）非空。对 \(i\) 用递降归纳法，我们断定 \(\mathbf{K}_0\) 非空，推论得证。

